% ECONOMICS_PROBLEM: {"id": "D72-413", "title": "Media markets, social media and polarization", "jel_code": "D72", "source_id": "OP-0123"}
% !TeX program = xelatex
\documentclass[11pt,a4paper]{article}
\usepackage[margin=23mm]{geometry}
\usepackage{fontspec}
\setmainfont{Latin Modern Roman}
\usepackage{amsmath,amssymb,mathtools}
\usepackage{xcolor,enumitem,booktabs,longtable,array,xurl,hyperref}
\definecolor{muted}{HTML}{53636C}
\hypersetup{colorlinks=true,urlcolor=blue,linkcolor=blue}
\setlength{\parindent}{0pt}
\setlength{\parskip}{5pt}
\newcommand{\R}{\mathbb R}
\newcommand{\E}{\mathbb E}
\newcommand{\N}{\mathbb N}
\newcommand{\ind}{\mathbf 1}
\newcommand{\Var}{\operatorname{Var}}
\newcommand{\Cov}{\operatorname{Cov}}
\newcommand{\supp}{\operatorname{supp}}
\DeclareMathOperator*{\argmin}{arg\,min}
\DeclareMathOperator*{\argmax}{arg\,max}
\newcommand{\entrymeta}[3]{{\sffamily\small\color{muted}\textbf{\texttt{#1}}\quad|\quad #2\par #3\par}}
\newcommand{\Problemref}[1]{\texttt{#1}}
\newcommand{\JELref}[2]{\texttt{#2} (JEL \texttt{#1})}
\begin{document}
\section*{Media markets, social media and polarization}
\textbf{Problem ID:} \texttt{D72-413}\quad \textbf{JEL:} \texttt{D72}\par
% BEGIN ECONOMICS STATEMENT
\label{problem:OP-0123}
\label{atlas:prob:Q3}
{\sffamily\small\color{muted}\textbf{Q3} | Political economy | Editorial research problem family\par}
\subsection*{Definitions and assumptions}
Fix a population $N=\{1,\ldots,n\}$ with $n\geq2$, a horizon $h$, and a media policy $a$ specifying channel availability, recommendation rules, content moderation and access costs. Let $X_{ih}(a)\in[-1,1]$ be a validated ideological score, $B_{ih}(a)$ a specified belief score, and $V_{ih}(a)$ voting behavior. Responses may depend on the entire population's exposure and on endogenous content supply. Define the population polarization index
\[
\Pi_h(a)=\frac{1}{n(n-1)}\sum_{i\ne j}(X_{ih}(a)-X_{jh}(a))^2.
\]
Fix the population composition or explicitly include migration and user entry. A model class $\mathcal M$ specifies exposure, selection, network interference, content supply and measurement. For an individual-assignment design, $z=(z_1,\ldots,z_n)\in\{0,1\}^n$ is the exposure vector under a fixed background policy; $X_{ih}(z)$ is the corresponding potential score. The notation $X_{ih}(z_i,z_{-i})$ uses this same map.
\subsection*{Formal research question}
Identify or sharply bound $\mathbb E[\Pi_h(a)-\Pi_h(a_0)]$ and the population mean changes in $B$ and $V$. For randomized individual exposure, additionally distinguish a direct contrast at fixed peer exposures from the total policy contrast:
\[
\delta_i(z_{-i})=\mathbb E[X_{ih}(1,z_{-i})-X_{ih}(0,z_{-i})].
\]
Determine what assignment designs or structural restrictions connect these individual contrasts to the population-level index after platform and political-supply responses. Give valid inference under the declared network dependence and sampling scheme. Identification of the polarization contrast requires joint or pairwise potential-score information, not just marginal mean effects.
\subsection*{Interpretation and scope}
The squared-distance index is one operational notion of ideological polarization, not a measure of every form of political conflict. Affective hostility, ideological extremity and electoral sorting require distinct outcomes. Randomized deactivation estimates effects of that intervention, not every algorithm change. Population-wide intervention effects include user substitution between media and strategic political messaging. Homophily and algorithmic amplification may be observationally equivalent without additional exposure variation.
\subsection*{Source audit and status}
All three atlas citations are verified, with publication and working-paper dates distinguished. The television papers provide setting-specific persuasion evidence; Allcott--Gentzkow provides fake-news exposure and belief evidence. They do not establish the displayed population-wide feed counterfactual or a general long-run polarization effect. The proposed formalization is an editorial research agenda whose application must state a concrete design.
\subsection*{References}
{\footnotesize\raggedright
\begin{enumerate}[label={\textbf{[S\arabic*]}},leftmargin=3em]
\item Stefano DellaVigna and Ethan Kaplan (2007). \emph{The Fox News Effect: Media Bias and Voting}. Quarterly Journal of Economics 122(3), 1187--1234.\par
\textit{Locator and source role:} Abstract and cable-channel introduction design; effects on voting, not a general estimate of polarization.\par
\url{https://academic.oup.com/qje/article-abstract/122/3/1187/1879517}
\item Ruben Enikolopov, Maria Petrova, and Ekaterina Zhuravskaya (2011). \emph{Media and Political Persuasion: Evidence from Russia}. American Economic Review 101(7), 3253--3285.\par
\textit{Locator and source role:} Abstract and media-exposure identification; setting-specific political persuasion.\par
\url{https://www.aeaweb.org/articles?id=10.1257/aer.101.7.3253}
\item Hunt Allcott and Matthew Gentzkow (2017). \emph{Social Media and Fake News in the 2016 Election}. Journal of Economic Perspectives 31(2), 211--236.\par
\textit{Locator and source role:} Abstract and browsing/survey evidence; descriptive exposure and belief evidence does not identify long-run feed-policy effects.\par
\url{https://www.aeaweb.org/articles?id=10.1257/jep.31.2.211}
\end{enumerate}
}

\subsection*{Common conventions: Source mathematical conventions}

Unless an entry states otherwise, agent and firm sets are finite. State and action spaces are measurable, strategies and outcome maps are measurable, probability laws are specified, and expectations exist for the targets under discussion. Infinite-horizon discount factors lie in $(0,1)$ and discounted objectives are finite. Norms, tolerances and complexity input sizes must be specified for computational comparisons. Constants or primitives that appear locally are inputs, not empirical facts asserted by this catalogue.

\textbf{Identification.} A statistical model $m\in\mathcal M$ implies an observable law $P_m$ and target $\theta(m)$. At law $P$, its identified set is
\[
\Theta(P;\mathcal M)=\{\theta(m):m\in\mathcal M,\ P_m=P\}.
\]
Point identification holds when this set is a singleton, or equivalently when $P_{m_1}=P_{m_2}$ implies $\theta(m_1)=\theta(m_2)$ for all compatible models. Sharp bounds mean bounds attained, or approached when the boundary is not attained, by compatible models. Writing this set is a definition, not a solution: the substantive task is to establish informative restrictions and derive or estimate the set. An unrestricted-confounding model may give only trivial bounds.

\textbf{Inference.} Identification, estimability and valid uncertainty quantification are separate questions. Fix $\alpha\in(0,1)$. A confidence procedure $C_n$ over a declared law class $\mathcal P$ has asymptotic uniform coverage at level $1-\alpha$ when
\[
\liminf_{n\to\infty}\inf_{P\in\mathcal P}\Pr_P\{\theta(P)\in C_n\}\geq1-\alpha.
\]
For set-identified targets the coverage event must be defined separately, for example coverage of the full identified set. If an entry uses identification-indexed models instead of uniquely indexed laws, the same coverage convention is applied over those models. Pointwise limits do not establish uniform validity, and asymptotic validity does not guarantee finite-sample precision. Sample sizes, dependence structures and weak-identification sequences are part of each application.

\textbf{Counterfactuals.} $Y_i(a)$ is a potential outcome under a specified intervention, including its timing, implementation and versions. It is not $Y_i$ conditional on voluntarily choosing $a$. A full intervention vector permits interference; shorthand scalar treatment notation does not silently impose no interference. Assignment, selection, attrition, anticipation and measurement must be specified. A claim about an instrument's local effect is not automatically a population policy effect.

\textbf{Economic equilibrium.} A counterfactual model includes best responses, constraints, feasibility, market clearing, state transitions and expectations consistent with the stated information. When several equilibria exist, either a declared selector is an input or the target is set-valued. Comparisons across policy regimes must use the stated selection convention. Reduced-form relationships are not presumed invariant after an equilibrium-changing intervention.

\textbf{Optimization and welfare.} Optimization is over the stated feasible set. If attainment is not established, $\sup$ or $\inf$ and $\varepsilon$-optimal choices replace an asserted maximizer or minimizer. For example, with value $V=\sup_{a\in\mathcal A}W(a)<\infty$, an $\varepsilon$-optimal set is $\{a\in\mathcal A:W(a)\geq V-\varepsilon\}$ for $\varepsilon>0$. Benefits, costs, risk and health or environmental outcomes can be aggregated only using declared units and conversion weights. Interpersonal utility weights, the treatment of fiscal transfers and foreign profits, discounting, and the behavioral welfare criterion are explicit inputs. A comparative-static derivative additionally requires existence and appropriate differentiability of the selected counterfactual.

\textbf{Sources.} Local labels [S1], [S2], and so forth restart in every question. Locators identify the relevant abstract, model, section or result at the level actually checked; a bibliographic or abstract-level verification is not represented as inspection of an entire proof. Working papers are distinguished from later publications and changing versions. Source summaries are paraphrased. All URL checks and editorial formulations refer to the audit date stated above.
% END ECONOMICS STATEMENT
\end{document}
